\chapter*{Abstract}
\addcontentsline{toc}{chapter}{Abstract}

Influence Diagrams are a standard formalism for decision-making under
uncertainty, but specifying their numerical parameters, the conditional
probability tables and the utilities, demands a costly and exhaustive elicitation
from domain experts, which is the main obstacle to their practical use. 

This
thesis investigates whether those parameters can instead be recovered
automatically from a small set of expert decision rules, assuming the structure
of the diagram is already known. The task is formulated as a constrained
black-box optimisation and addressed with Estimation of Distribution Algorithms:
the optimiser searches over an unconstrained real vector that a decoder always
maps to a valid model, and each candidate is scored by how faithfully the policy
it induces reproduces the given rules. A proof-of-concept system was built on the
\textsc{SMILE} and \textsc{EDAspy} libraries and evaluated on two clinical
diagrams of different size, comparing five loss functions and four EDA variants.


The results show that recovery is feasible and improves as more rules are added,
but that reproducing the training rules does not guarantee recovering the full
policy: the problem is degenerate. The choice of optimiser is the dominant
factor, and expressive variants that scale to the larger diagram are the
limiting resource. The information carried by the rules resides in the utilities
rather than in the probabilities, so recovering only the utilities is both
cheaper and more accurate. Finally, the composition of the rule set matters as
much as its size, and the most informative rules are the decisive ones. The work
provides a working recovery method, a characterisation of its limits, and a set
of concrete directions for future research.

\vspace{1em}
\noindent\textbf{Keywords:} Influence Diagrams; parameter elicitation; decision
rules; Estimation of Distribution Algorithms; black-box optimisation; decision
support systems; utility recovery.



%%--------------
\newpage
%%--------------

\chapter*{Resumen}
\addcontentsline{toc}{chapter}{Resumen}

Los \emph{Diagramas de Influencia} (DIs) son un formalismo gráfico compacto para
la toma de decisiones bajo incertidumbre: un único grafo codifica las variables
inciertas, las opciones disponibles para el agente y las preferencias sobre los
resultados. Su principal obstáculo práctico no es el grafo, sino sus números.
Rellenar cada tabla de probabilidad condicional (TPC) y cada valor de utilidad
exige largas entrevistas con expertos escasos, y el esfuerzo crece tan rápido
con el tamaño del modelo que se convierte en el coste dominante de desplegar un
DI en dominios reales como la ayuda a la decisión clínica.

Esta tesis estudia una alternativa a esta elicitación exhaustiva. Suponemos que
la estructura del DI ya está dada y nos preguntamos si sus parámetros numéricos
pueden \emph{recuperarse automáticamente} a partir de un pequeño conjunto de
\emph{reglas de decisión} de la forma ``si se cumplen estas condiciones, toma
esta acción''. Las reglas están mucho más cerca de la forma en que razonan los
expertos que las probabilidades en bruto, y por lo general solo se necesitan unas
pocas para describir una política. El problema se plantea como una optimización
de caja negra con restricciones sobre el espacio conjunto de TPCs y utilidades,
y se resuelve con \emph{Algoritmos de Estimación de Distribuciones} (EDAs), una
familia de optimizadores evolutivos que guían la búsqueda con un modelo
probabilístico de las mejores soluciones encontradas hasta el momento.

El proceso de recuperación se implementó sobre el motor de inferencia
\textsc{SMILE} y la biblioteca \textsc{EDAspy}. El optimizador busca sobre un
vector real sin restricciones que un decodificador convierte siempre en un
modelo válido (softmax para las tablas de probabilidad, sigmoide para las
utilidades), de modo que ningún candidato tiene que repararse ni descartarse.
Cada candidato se puntúa según lo bien que la política que induce reproduce las
reglas del experto, empleando una de cinco funciones de pérdida, y se comparan
cuatro variantes de EDA de expresividad creciente (UMDA, EMNA, EGNA y KEDA). El
método se evalúa sobre dos diagramas clínicos de distinto tamaño: el pequeño
\emph{bypass} (manejo de la cardiopatía coronaria) y el mayor \emph{NHLy}
(manejo del linfoma no Hodgkin).

El estudio experimental arroja cinco hallazgos principales. Primero, la
recuperación es \emph{factible}: mostrar solo una fracción de las reglas eleva la
política recuperada muy por encima del azar, y la precisión crece a medida que se
añaden más reglas. Segundo, y de forma central, reproducir las reglas de
entrenamiento \emph{no} es lo mismo que recuperar la política: muchos vectores de
parámetros satisfacen las reglas mostradas pero discrepan en las no vistas, por
lo que el problema está parcialmente condicionado. Tercero, el \emph{optimizador}
es el factor dominante: KEDA es el más potente en la red pequeña pero escala mal,
de modo que en el diagrama mayor solo los más económicos EMNA y UMDA resultan
asequibles, con EMNA claramente por delante. Cuarto, lo que las reglas realmente
determinan son las \emph{utilidades}, no las probabilidades: buscar solo las
utilidades reproduce la política casi a la perfección, mientras que añadir las
tablas de probabilidad a la búsqueda resulta incluso perjudicial. Quinto,
\emph{qué} reglas se muestran importa tanto como cuántas: para un número fijo de
reglas, la elección genera una amplia dispersión en la precisión, y las reglas
más informativas son las decisivas, aquellas con la mayor diferencia de utilidad
entre la acción recomendada y su alternativa.

En resumen, el cuello de botella del método es la degeneración del problema y la
expresividad de la búsqueda. La receta práctica que se deriva consiste en recuperar las
utilidades manteniendo fijas las probabilidades, combinar un optimizador
expresivo con una pérdida sencilla e invertir el esfuerzo del experto en un
pequeño conjunto de reglas decisivas. El trabajo entrega una prueba de concepto
validada y señala la selección activa de reglas, la explotación de la población
degenerada y el escalado a diagramas mayores como las continuaciones más
prometedoras.


\newpage
